\documentclass{article}

\usepackage{amsmath,amssymb}
\usepackage{xurl}
\usepackage[colorlinks=true,linkcolor=blue,citecolor=blue,urlcolor=blue]{hyperref}

% Shared commands for notes in docs/paper-gaps/.

\DeclareUnicodeCharacter{2080}{\ifmmode{}_{0}\else\textsubscript{0}\fi}
\DeclareUnicodeCharacter{2081}{\ifmmode{}_{1}\else\textsubscript{1}\fi}
\DeclareUnicodeCharacter{2082}{\ifmmode{}_{2}\else\textsubscript{2}\fi}
\DeclareUnicodeCharacter{2099}{\ifmmode{}_{n}\else\textsubscript{n}\fi}

\newcommand{\Lean}{\textsc{Lean}}
\ExplSyntaxOn
\NewDocumentCommand{\leanid}{m}{
  \ifmmode
    \textsf{\detokenize{#1}}
  \else
    \group_begin:
    \urlstyle{sf}
    \tl_set:Nn \l_tmpa_tl {#1}
    \tl_replace_all:Nnn \l_tmpa_tl {\_} {_}
    \exp_args:NV \path \l_tmpa_tl
    \group_end:
  \fi
}
\ExplSyntaxOff
\providecommand{\tr}{\operatorname{tr}}
\newcommand{\tnleanIssueBase}{https://github.com/LionSR/TNLean/issues/}
\newcommand{\tnleanPrBase}{https://github.com/LionSR/TNLean/pull/}
\newcommand{\tnleanDocBase}{https://sirui-lu.com/TNLean/docs/}
\newcommand{\ghissue}[1]{\href{\tnleanIssueBase#1}{GitHub issue \##1}}
\newcommand{\ghpr}[1]{\href{\tnleanPrBase#1}{PR \##1}}
\newcommand{\doclink}[1]{\href{\tnleanDocBase#1.html}{\path{#1}}}

% Dirac notation and the identity operator, as in blueprint/src/macros/common.tex.
% \providecommand keeps note-local definitions authoritative.
\usepackage{dsfont}
\providecommand{\ket}[1]{|#1\rangle}
\providecommand{\bra}[1]{\langle #1|}
\providecommand{\braket}[2]{\langle #1 | #2 \rangle}
\providecommand{\ketbra}[2]{|#1\rangle\!\langle #2|}
\providecommand{\Id}{\mathds{1}}
\providecommand{\one}{\mathds{1}}
\providecommand{\C}{\mathbb{C}}
\providecommand{\MN}[1]{M_{#1}(\C)}
\providecommand{\MD}{M_D(\C)}

% External referencing for paper sources.  The build helper writes sanitized
% auxiliary files under build/paper-aux/ and excludes duplicated source labels.
\usepackage{xr}

% Import a generated paper auxiliary file with a caller-supplied label prefix.
% The candidate paths support builds from docs/paper-gaps/, the repository root,
% and build/paper-aux/ itself.
\newcommand{\importpaperaux}[2]{%
  \IfFileExists{../../build/paper-aux/#2.aux}{%
    \externaldocument[#1]{../../build/paper-aux/#2}%
  }{%
    \IfFileExists{build/paper-aux/#2.aux}{%
      \externaldocument[#1]{build/paper-aux/#2}%
    }{%
      \IfFileExists{#2.aux}{%
        \externaldocument[#1]{#2}%
      }{}%
    }%
  }%
}

% General external reference macros
\newcommand{\paperref}[2]{\ref{#1-#2}}
\newcommand{\papereqref}[2]{(\ref{#1-#2})}

% CPSV16 source (1606.00608 / MPDO-22-12-17-2.tex) external document and shortcuts
\importpaperaux{cpsv16-}{1606.00608/MPDO-22-12-17-2-tnlean}
\newcommand{\cpsvref}[1]{\paperref{cpsv16}{#1}}
\newcommand{\cpsveqref}[1]{\papereqref{cpsv16}{#1}}

% Verdict marker: \gapnote{<kind>}{<status>}. Kind is the policy
% classification (clarification, local-correction, scope-restriction,
% unfaithful, false-source, open-gap); status is open, wip (elimination
% underway), resolved, or historical. Machine-read by the site index;
% prints a verdict line.
\newcommand{\gapnote}[2]{\par\noindent\textbf{Verdict:} #1 (#2).\par}


\title{Power Sums in the BNT Weight Comparison}
\author{}
\date{2026-05-08}

\begin{document}
\maketitle
\gapnote{clarification}{historical}

\section{The source lemma}

This note concerns the scalar power-sum step in the non-periodic Fundamental
Theorem proof for matrix product states in canonical form. The source is
Lemma~\texttt{Lem:app\_simple} of
Ref.~\cite[lines~1155--1163]{Cirac2016MPDO_arXiv}. It considers two finite
families
\begin{align}
    (\lambda_{a,k})_{k=1}^{x_a},
    \qquad
    (\lambda_{b,k})_{k=1}^{x_b},
    \label{eq:cpsv16_power_sum_alternative_route_source_families}
\end{align}
ordered by modulus and phase. Its hypothesis, printed as
Equation~\texttt{eq:app\_aux} of Ref.~\cite{Cirac2016MPDO_arXiv}, is
\begin{align}
    \sum_{k=1}^{x_a}\lambda_{a,k}^{N}
        &=\sum_{k=1}^{x_b}\lambda_{b,k}^{N},
        \qquad
        1\leq N\leq\max\{x_a,x_b\}.
    \label{eq:cpsv16_power_sum_alternative_route_source_power_sums}
\end{align}
The conclusion is
\begin{align}
    x_a
        &=x_b,
    \label{eq:cpsv16_power_sum_alternative_route_source_cardinality}\\
    \lambda_{a,k}
        &=\lambda_{b,k},
        \qquad 1\leq k\leq x_a.
    \label{eq:cpsv16_power_sum_alternative_route_source_ordered_equality}
\end{align}

In the equal-case corollary of the Fundamental Theorem in
Ref.~\cite{Cirac2016MPDO_arXiv}, this lemma is applied after matching the BNT
representatives. For a fixed representative $j$, the coefficient families are
\begin{align}
    (\mu_{j,q})_{q=1}^{r_{a,j}},
    \qquad
    (\nu_{j,q}e^{i\phi_j})_{q=1}^{r_{b,j}}.
    \label{eq:cpsv16_power_sum_alternative_route_bnt_weight_families}
\end{align}
The formal statements replace the source ordering convention with equality of
multisets. They retain the nonzero-coefficient hypothesis because positive
powers cannot count zero entries.

\section{The source-facing formal route}

Fix a BNT representative $j$. BNT linear independence first gives eventual
equality of the power sums:
\begin{align}
    \sum_{q=1}^{r_j^S}(\mu^S_{j,q})^N
        &=\sum_{q=1}^{r_j^T}(\mu^T_{j,q})^N.
    \label{eq:cpsv16_power_sum_alternative_route_eventual_power_sum_equality}
\end{align}
A finite signed sum of nonzero geometric sequences that vanishes eventually
vanishes at every exponent. Hence
(\ref{eq:cpsv16_power_sum_alternative_route_eventual_power_sum_equality}) holds
for every positive $N$.

The formal proof then applies the unequal-cardinality finite-range power-sum
theorem.\footnote{Formal declaration:
\leanid{Matrix.sum_pow_eq_implies_card_eq_and_multiset_eq_of_le_max_card}.}
The BNT comparison uses the corresponding sector-weight lemma.\footnote{Formal
declaration: \leanid{MPSTensor.matched_sector_weight_multiset_eq} in
\path{TNLean/MPS/FundamentalTheorem/SectorBNT/WeightEquiv.lean}. It composes
\leanid{MPSTensor.SectorWeightData.power_sums_eq_of_eventually_eq_hetero}
with the unequal-cardinality power-sum theorem.}
Its conclusion is
\begin{align}
    r_j^S
        &=r_j^T,
    \label{eq:cpsv16_power_sum_alternative_route_formal_cardinality}\\
    \{\!\{\mu^S_{j,q}:1\leq q\leq r_j^S\}\!\}
        &=\{\!\{\mu^T_{j,q}:1\leq q\leq r_j^T\}\!\}.
    \label{eq:cpsv16_power_sum_alternative_route_formal_multiset}
\end{align}
This is the multiset form of
(\ref{eq:cpsv16_power_sum_alternative_route_source_ordered_equality}).

\section{The alternative formal route}

An earlier formal proof reached the same endpoint by using the recurrence for
finite sums of geometric sequences before applying Newton--Girard. Let
\begin{align}
    F(N)
        &=\sum_{\ell=1}^{m}c_\ell w_\ell^N,
        \qquad w_\ell\neq 0.
    \label{eq:cpsv16_power_sum_alternative_route_geometric_sum}
\end{align}
Assume that $F(N)=0$ for every $N\geq M$. Subtracting $w_1F(N)$ from
$F(N+1)$ gives
\begin{align}
    F(N+1)-w_1F(N)
        &=\sum_{\ell=1}^{m}c_\ell w_\ell^{N+1}
          -\sum_{\ell=1}^{m}c_\ell w_1w_\ell^N
    \label{eq:cpsv16_power_sum_alternative_route_recurrence_expansion}\\
        &=\sum_{\ell=2}^{m}
          c_\ell(w_\ell-w_1)w_\ell^N
    \label{eq:cpsv16_power_sum_alternative_route_reduced_sum}\\
        &=0,
        \qquad N\geq M.
    \label{eq:cpsv16_power_sum_alternative_route_reduced_sum_eventual_zero}
\end{align}
The sum in
(\ref{eq:cpsv16_power_sum_alternative_route_reduced_sum}) has one fewer term.
Induction on $m$ shows that it vanishes for every $N\geq 0$. Therefore,
\begin{align}
    F(N+1)
        &=w_1F(N),
        \qquad N\geq 0.
    \label{eq:cpsv16_power_sum_alternative_route_global_recurrence}
\end{align}
Iteration gives
\begin{align}
    F(k)
        &=w_1^kF(0),
        \qquad k\geq 0.
    \label{eq:cpsv16_power_sum_alternative_route_recurrence_solution}
\end{align}
Since $F(M)=0$ and $w_1^M\neq 0$,
(\ref{eq:cpsv16_power_sum_alternative_route_recurrence_solution}) implies
$F(0)=0$. Thus $F(k)=0$ for every $k\geq 0$.

For the two sector-weight families, apply this lemma to
\begin{align}
    F(N)
        &=\sum_{q=1}^{r_j^S}(\mu^S_{j,q})^N
          -\sum_{q=1}^{r_j^T}(\mu^T_{j,q})^N.
    \label{eq:cpsv16_power_sum_alternative_route_signed_weight_sum}
\end{align}
The nonzero-weight hypothesis permits the use of $w_1^M\neq 0$. Every weight
in \leanid{SectorWeightData} is nonzero, so the exponent-zero power sum counts
the entries:
\begin{align}
    \sum_{q=1}^{r_j}\mu_{j,q}^{0}
        &=r_j.
    \label{eq:cpsv16_power_sum_alternative_route_zero_power_count}
\end{align}
Hence $F(0)=0$ gives
(\ref{eq:cpsv16_power_sum_alternative_route_formal_cardinality}). After
transporting the right-hand index set along this equality, the extrapolation
also gives
\begin{align}
    \sum_{q=1}^{r_j^S}(\mu^S_{j,q})^N
        &=\sum_{q=1}^{r_j^S}(\mu^T_{j,q})^N,
        \qquad N\geq 1.
    \label{eq:cpsv16_power_sum_alternative_route_common_index_power_sums}
\end{align}
The same-cardinality Newton--Girard theorem then recovers the weight multiset.

This route previously existed as checked \Lean{} code through declarations for
the copy count, transported positive power sums, and weight-multiset
equality.\footnote{Formal declarations:
\leanid{MPSTensor.SectorWeightData.copies_eq_of_eventually_coeff_eq};
\leanid{MPSTensor.SectorWeightData.eventually_coeff_eq_implies_all_pos_eq_after_copies};
and
\leanid{MPSTensor.SectorWeightData.weight_multiset_eq_of_copies_eq_of_coeff_eq}.}
Those declarations have been removed because the proportional Fundamental
Theorem no longer uses this same-cardinality detour. The calculation records
the removed route for comparison with Ref.~\cite{Cirac2016MPDO_arXiv}; it is
not a present \Lean{} proof.

\section{Conclusion}

Under the former nonzero-weight hypotheses, the removed route has no
mathematical gap. The extrapolation argument proves the $N=0$ identity, and
(\ref{eq:cpsv16_power_sum_alternative_route_zero_power_count}) justifies the
copy count. The source-facing theorem surface should use the
unequal-cardinality power-sum lemma because that is the lemma invoked after
BNT representative matching in Ref.~\cite{Cirac2016MPDO_arXiv}.

Without a nonzero-coefficient hypothesis, positive power sums cannot count
zero weights. Adjoining a zero to a finite family leaves every positive power
sum unchanged. A theorem without this hypothesis can therefore recover only
the multiset obtained after deleting zero entries, unless it separately
assumes an exponent-zero identity.

\bibliographystyle{alpha}
\bibliography{references}

\end{document}
